% =============================================================================
% CASPEr@night --- Agilent/VnmrJ Quickstart (SIU Carbondale)
% Single-recipient quickstart for Agilent sessions on the partner instrument
% (Boyd Goodson's 400 MHz DD2, VnmrJ 3.2, 5 mm HCN, room temperature);
% two sessions ran on it 2026-09-14.
% v0.7.1: randomized repeated gain ladder (2.2), spin-noise tuning ladder
% (2.6), axis-sign check (2.7), physical probe temperatures (2.1, 3.2);
% roles now come from the save names and rg_ladder from the rungs.
% Distilled from vendors/agilent/README.md (Tier 1 + partner checklist),
% packer/answers.example.json, and uploader/upload_bundle.py.
% Build: pdflatex x2.
% =============================================================================
\documentclass[11pt]{article}

\usepackage[T1]{fontenc}
\usepackage[utf8]{inputenc}
\usepackage{lmodern}
\usepackage{microtype}
\usepackage[margin=2.4cm]{geometry}
\usepackage{xcolor}
\usepackage{booktabs}
\usepackage{enumitem}
\usepackage{listings}
\usepackage{tcolorbox}
\usepackage[hidelinks]{hyperref}

\definecolor{cmdbg}{gray}{0.955}
\definecolor{cmdframe}{gray}{0.70}
\definecolor{boxframe}{rgb}{0.15,0.30,0.45}
\definecolor{boxbg}{rgb}{0.945,0.965,0.98}

\lstset{
  basicstyle=\ttfamily\small,
  backgroundcolor=\color{cmdbg},
  frame=single,
  rulecolor=\color{cmdframe},
  framesep=4pt,
  xleftmargin=6pt,
  xrightmargin=6pt,
  breaklines=true,
  breakatwhitespace=true,
  postbreak=\mbox{\space\space},
  columns=fullflexible,
  keepspaces=true,
  aboveskip=4pt,
  belowskip=4pt,
  literate={~}{\textasciitilde}{1}
}

\newcommand{\hnmr}{\textsuperscript{1}H}
\newcommand{\dtwo}{D\textsubscript{2}O}
\newcommand{\htwo}{H\textsubscript{2}O}
\newcommand{\file}[1]{\texttt{#1}}

\setlist[itemize]{topsep=3pt, itemsep=2pt, parsep=0pt}
\setlist[enumerate]{topsep=3pt, itemsep=3pt, parsep=0pt}

\emergencystretch=1.5em

\title{\textbf{CASPEr@night --- Agilent/VnmrJ Quickstart}\\
\large (SIU Carbondale --- 400\,MHz DD2, VnmrJ 3.2)}
\author{John W. Blanchard (\href{mailto:jwbquantum@gmail.com}{jwbquantum@gmail.com})
\and Reza Ebadi \and Claude (Anthropic)}
\date{Software v0.7.1 --- September 16, 2026}

\begin{document}
\maketitle

\section{What this is}

This guide covers one specific session: a receive-only spin-noise
measurement on the SIU Carbondale 400\,MHz Agilent DD2 (VnmrJ~3.2, 5\,mm
HCN probe, room temperature). It is the Agilent leg of the CASPEr@night
network --- the same physics protocol the Bruker fleet runs, recording
the spin-noise feature of water on a room-temperature probe (the
Gu\'eron absorption dip in equilibrium; the September sessions here
showed a +27\% emission bump instead, which is why the tuning ladder of
2.6 now exists) ---
acquired by hand with ordinary VnmrJ tools (\texttt{s2pul},
\texttt{go}, \texttt{svf}) and converted into a network bundle
afterward on any laptop. Nothing in the session pulses the sample above
tune-and-match power: the calibration steps use ${\sim}1^{\circ}$ tips
at your ordinary observe power, and the noise blocks use
\texttt{pw=0} at the bottom of the transmitter power scale. One thing
to know going in: \textbf{this path has run on real hardware on one
instrument --- this one.} Two sessions on 2026-09-14 (68 and 66
experiments) were packed, validated and uploaded with zero reader
warnings, so the converter has parsed real files from one instrument of
this console family and no other. If your files make it choke, that is still data we need ---
ugly output is useful output, and Section~\ref{sec:wrong} says what to
send. You never need the GitHub repository, git, or a README: this PDF
plus the \file{config.json} attached to the same e-mail is everything.

\begin{tcolorbox}[colback=boxbg, colframe=boxframe,
  title=\textbf{What you need}, fonttitle=\bfseries]
\begin{enumerate}[leftmargin=1.4em]
\item \textbf{Any laptop or desktop with Python 3.6 or newer} ---
      macOS, Linux, and Windows all work; nothing is installed on the
      spectrometer. Check with:
\begin{lstlisting}
python3 --version
\end{lstlisting}
\item \textbf{The software --- one download, no git.} Open
\begin{lstlisting}
https://github.com/CASPEr-Night/spin-noise-network/releases/latest
\end{lstlisting}
      and under \emph{Assets} click \textbf{Source code (zip)}. Unzip
      it anywhere.
\item \textbf{The \file{config.json} from John's e-mail} --- drop it
      into the unzipped folder at exactly
      \file{uploader/config.json} (next to
      \file{upload\_bundle.py}). It is your facility's upload key ---
      treat it like a password.
\item \textbf{A 5\,mm tube of 90\% \htwo{} / 10\% \dtwo, lightly
      doped} --- a retired shim/lineshape sample is fine, as long as
      you know (or can best-guess) the composition. The \htwo{}
      fraction \emph{is} the measurement.
\item \textbf{About 1\,h\,20\,min of instrument time} --- roughly
      20\,min of setup and calibration you attend, ${\geq}30$\,min of
      noise blocks the console mostly runs itself, then some 15\,min
      for the tuning ladder and the axis-sign check (2.6--2.7; about
      25\,min with the six-block ladder a 1\%-class tube needs).
\end{enumerate}
\end{tcolorbox}

\section{At the spectrometer}
\label{sec:spectrometer}

Everything below is standard VnmrJ operation --- no macro, nothing
installed on the console. Save every experiment with \texttt{svf} into
one session directory, one \file{.fid} per step; the leading number in
each save name identifies the experiment and the \texttt{\_sn\_\ldots}
part tells the converter its role, so keep both as printed. Anything
after them is free as long as it does not run straight on in letters:
\texttt{17\_sn\_noise\_b} and \texttt{17\_sn\_noise2} are noise blocks,
\texttt{17\_sn\_noisetest} is not recognised and the converter stops
and says so.

\subsection*{2.1\quad Setup (nothing new here, ${\sim}15$\,min)}

Insert the water tube, set or record the temperature (298\,K if VT is
off and the bore is at room temperature), and let it equilibrate a few
minutes. Tune and match \hnmr{} by your normal manual wobble procedure
--- no autotune on this probe, and none needed. Shim as usual, with
one caution. On a near-neat water tube the water line may sit at
${\sim}30$\,Hz wide and asymmetric, and gradient shimming cannot
narrow it --- that is radiation damping, the probe's back-action on
the spins, not a shim fault. On this instrument, identical shims and
identical acquisition gave an 8.5\,Hz symmetric line on a 1\%
\htwo{}/\dtwo{} Gd-doped tube and a 32\,Hz asymmetric line on the
90/10 tube; three gradient-shim iterations on the 90/10 tube moved
z3--z5 by several hundred DAC units each and changed the linewidth by
under 0.2\,Hz. Judge shim quality on the lock signal or on a
dilute/doped tube, not on the near-neat water line, and do not judge a
shim set by the size of its DAC values --- the set that gave the
8.5\,Hz line had z3 and z3y past $\pm$8000. The broad line is the
regime that produces a large noise feature, so leave it as it is.

Lock: run with the lock \textbf{on} during the noise blocks unless you
have a reason not to. On this instrument the locked session drifted
0.25\,Hz between opening and closing reference over a 38\,min bracket,
with an unchanged lineshape; the unlocked session drifted 0.56\,Hz over
a 78\,min bracket. That is 0.0065 versus 0.0072\,Hz/min --- the same
drift rate within what two brackets can show --- so the two sessions do
not by themselves measure what the lock buys. (\texttt{alock='n'} in
the block parameters below only disables automatic re-locking; it does
not release an established lock.) Either way, \textbf{record the state}
--- the questionnaire asks for it. Calibrate the \hnmr{} 90$^{\circ}$
pulse or take the probe-file value, and \textbf{write down pw90 and
tpwr} --- the questionnaire wants both, and the closing reference
must reuse them. Start from your standard \hnmr{} water setup in
\texttt{s2pul} and leave \texttt{sw} at your usual value throughout.

\textbf{Probe temperatures.} The questionnaire (3.2) asks for the coil
and preamplifier temperatures, and it means \emph{physical}
temperatures: on a room-temperature probe the coil sits at the
probe-body temperature, which is the lab ambient, and the preamplifier
at its own physical ambient. Read the lab thermometer (or the bore/VT
readout with VT off) while you are at the magnet and enter that --- a
measured 296.5\,K beats an assumed 298\,K; \texttt{null} is for having
no reading at all. If a better number turns up after the upload, e-mail
it to the maintainer: corrections are applied at analysis time, as
numbers, never by re-uploading.

\subsection*{2.2\quad Gain ladder (${\sim}3$\,min)}

A randomized, repeated ladder of one-scan tiny-flip 1Ds --- the
receiver-linearity calibration. First find your ladder maximum: the
highest \texttt{gain} at which a tiny-flip 1D of this tube does not
overflow the receiver. \textbf{It is strongly sample-dependent} --- on
this probe 60\,dB was fine on a 1\% \htwo{} Gd-doped tube and the 90/10
\htwo{}/\dtwo{} tube needed 40\,dB. The ladder then visits every level
from 0 to that maximum in 5\,dB steps (legal gains are integer dB: 0,
5, \ldots, 40 for a 40\,dB maximum, nine levels; thirteen for 60\,dB),
\textbf{each level three times, in one random order}, with
\texttt{pad=1} so the receiver settles for a second after every gain
change. The September sessions used four ascending gains, and in both
the rungs' amplitude per unit gain fell by about 10\% from the bottom
rung to the top one (11\% and 10\%) --- which a monotonic ladder
cannot attribute: receiver
compression and signal drift are degenerate when gain and time rise
together. Visiting each level three times at random moments separates
them: compression follows the gain, drift follows the clock.
Twenty-seven acquisitions of a few seconds each, about 3 minutes in
all.

Set \texttt{pw} to pw90/90 --- 0.11\,$\mu$s for a 10\,$\mu$s pw90,
about a $1^{\circ}$ tip --- and \texttt{tpwr} to your calibrated
observe power (57 below stands in for yours). Write both values down:
the references in 2.3, 2.5 and 2.7 must use exactly these. Never use
\texttt{gain='n'} (autogain), and wait for each acquisition to finish
before saving. The console may store a different gain than you typed
--- asked for 13.3 and 26.7\,dB, this console stored 14 and 26 ---
which is why the converter reads the stored gain from each rung's
\file{procpar} and you list nothing. Save the rungs under four-digit
numbers counting up from 1000 in the order acquired (a number that can
never collide with a noise block); the gain in the name is for your
eyes only. One random order for a 40\,dB maximum, drawn once and
printed here so you need not draw it at the console --- any other
random order is as good, and the order that ran is on record in each
file's time stamp:

\begin{center}\small
\begin{tabular}{@{}l*{9}{r}@{}}
\toprule
save as & 1000 & 1001 & 1002 & 1003 & 1004 & 1005 & 1006 & 1007 & 1008\\
\texttt{gain} & 35 & 35 & 5 & 40 & 15 & 30 & 10 & 15 & 40\\
\midrule
save as & 1009 & 1010 & 1011 & 1012 & 1013 & 1014 & 1015 & 1016 & 1017\\
\texttt{gain} & 10 & 25 & 30 & 20 & 5 & 40 & 30 & 35 & 0\\
\midrule
save as & 1018 & 1019 & 1020 & 1021 & 1022 & 1023 & 1024 & 1025 & 1026\\
\texttt{gain} & 20 & 0 & 0 & 10 & 15 & 25 & 5 & 25 & 20\\
\bottomrule
\end{tabular}
\end{center}

The first three rungs, then, are:

\begin{lstlisting}
pw=0.11
tpwr=57
nt=1
pad=1
gain=35
go
svf('1000_sn_ladder_35')
gain=35
go
svf('1001_sn_ladder_35')
gain=5
go
svf('1002_sn_ladder_05')
\end{lstlisting}

and so on down the table. For a 60\,dB maximum use the levels 0, 5,
\ldots, 60 (13 levels, 39 rungs, about 4 minutes), each three times,
in an order you shuffle yourself before you start --- write it down
first. Read every \texttt{gain=40} below as your ladder-maximum value.

\subsection*{2.3\quad Opening reference (${\sim}1$\,min)}

One more tiny-flip 1D at the ladder-max gain, with a 2\,s record.
\texttt{pw} and \texttt{tpwr} are the two values you wrote down in
2.2 (the same stand-ins below), and \texttt{pad} goes back to 0 --- the
ladder left \texttt{pad=1} behind:

\begin{lstlisting}
pw=0.11
tpwr=57
pad=0
at=2
gain=40
go
svf('11_sn_ref_open')
\end{lstlisting}

\subsection*{2.4\quad Noise blocks (${\geq}30$\,min --- the point of the session)}

No pulse at all --- these lines are the receiver-only idiom of
Agilent's own \texttt{cryo\_noisetest} macro, with \texttt{tpwr=-16}
the bottom of the standard transmitter power scale:

\begin{lstlisting}
pw=0
tpwr=-16
nt=1
ss=0
pad=0
wshim='n'
alock='n'
gain=40
at=30
go
svf('12_sn_noise')
\end{lstlisting}

Then repeat until you have at least 30 minutes of noise data in total
--- about 60 blocks at \texttt{at=30} --- counting up from 17 (13 is
reserved for the closing reference):

\begin{lstlisting}
go
svf('17_sn_noise')
go
svf('18_sn_noise')
\end{lstlisting}

Three notes --- two settled on this console, one still open:
\begin{itemize}
\item \textbf{Record length.} The console caps a record at a point
      count, \texttt{np} $\leq$ 524288, so the longest record is
      $\texttt{at}_{\max} = 524288/(2\,\texttt{sw})$ --- about 41\,s
      at \texttt{sw} = 6410\,Hz. \texttt{at=30} with about 60 blocks
      gives 30 minutes; both sessions here used 30\,s records (62 and
      60 blocks). A narrower \texttt{sw} buys a proportionally longer
      record --- 3205\,Hz allows about 82\,s --- if your usual water
      window permits it.
\item \textbf{Transmitter silence.} Whether the DD2's transmitter
      chain is measurably silent at \texttt{pw=0 tpwr=-16} has not
      been checked --- no spectrum analyser was available on
      2026-09-14. If the noise floor shows something odd --- a
      carrier spike, a ridge --- do not chase it: finish the session
      and describe it in the questionnaire notes.
\item \textbf{Save names.} \texttt{svf} accepts names starting with a
      digit on this console --- both sessions saved every experiment
      under the \file{NN\_\ldots} pattern. The leading number and the
      \texttt{\_sn\_noise} part are parsed; anything after them is
      free unless it runs straight on in letters
      (\texttt{17\_sn\_noise\_b}, \texttt{17\_sn\_noise2} and
      \texttt{17\_sn\_noise-1} are read as noise blocks;
      \texttt{17\_sn\_noisetest} is not).
\end{itemize}

\subsection*{2.5\quad Closing reference (${\sim}1$\,min)}

Repeat 2.3 exactly --- same \texttt{at}, same \texttt{gain}, and
\texttt{pw} and \texttt{tpwr} back at the two values you wrote down in
2.2 (the noise blocks left \texttt{pw=0 tpwr=-16} behind, so both
must be restored; do not retype the stand-in numbers below unless they
are yours). The two references are a matched drift bracket, and a
closing reference at a different flip angle cannot be compared with
the opening one --- the first session here took them at
\texttt{pw}=0.05 and 0.1125\,$\mu$s, a $2.25\times$ mismatch in flip
angle, for exactly this reason. Save as 13:

\begin{lstlisting}
pw=0.11
tpwr=57
at=2
gain=40
go
svf('13_sn_ref_close')
\end{lstlisting}

That is the main session; 2.6 and 2.7 add two short calibration series
to it. An external acquisition driver --- a
Python~2.6 script that runs on the console workstation itself,
submits one block at a time to a running VnmrJ through
\texttt{listenon}/\texttt{send2Vnmr}, and waits for each \file{.fid}
to land before submitting the next --- ran the second session here end
to end on 2026-09-14 (66 of 66 blocks, unattended). It is being
contributed to \file{vendors/agilent/}; until it ships there, the
manual steps above are the recommended path. The draft VnmrJ macro
(\file{spin\_noise\_run.mac}) has still not run on hardware.

\subsection*{2.6\quad Spin-noise tuning ladder (15--25\,min)}

The sign and shape of the spin-noise feature depend on how the probe
is tuned. Equilibrium physics expects the Gu\'eron absorption dip on a
room-temperature probe; both September sessions on this probe showed a
clear feature of the opposite sign --- an emission bump, +27\% of the
floor on the 1\% \htwo{} tube --- which says the probe was tuned away
from the spin-noise tuning optimum, a setting that is in general not
the wobble minimum you tune to. Stepping the tuning changes the shape
of the feature as well as its size: at the optimum the dispersive
admixture vanishes and the feature is the purest absorptive dip; away
from it a dispersive component grows, and at the setting where the
feature nulls and flips from bump to dip the absorptive part is
passing through zero, so what remains is nearly pure dispersion ---
that crossover is a setting \emph{far} from the optimum, not the
optimum itself. The report fits every setting with the same
absorptive-plus-dispersive model as the main block, ranks the settings
that reach 5$\sigma$ by their dispersive fraction $|b/a|$, and names
the smallest as nearest the optimum; the null-or-flip setting is the
one it ranks farthest, so do not label it the optimum in your notes.
So after the closing reference, step the tuning
deliberately and record noise blocks at each setting: \textbf{at
least five settings bracketing your normal one} --- the normal tuning,
then one and two steps to either side, where a step is whatever you
can set and read reproducibly (a number of turns of the tune or match
capacitor, a reflected-power reading on the wobble display, a
wobble-curve minimum in dB). Same no-pulse parameters as 2.4,
\texttt{at=30}, two blocks per setting on the recommended 90/10 tube;
setting 0 is your normal tuning, saved as 2000 and 2001:

\begin{lstlisting}
pw=0
tpwr=-16
nt=1
ss=0
pad=0
wshim='n'
alock='n'
gain=40
at=30
go
svf('2000_sn_tune_0')
go
svf('2001_sn_tune_0')
\end{lstlisting}

Then retune to setting 1 and save \texttt{2010\_sn\_tune\_1} and
\texttt{2011\_sn\_tune\_1}, setting 2 as 2020/2021, and so on: setting
$k$ goes to 20$k$0 and 20$k$1, and the trailing \texttt{\_}$k$ tells
the converter which setting a block belongs to. \textbf{How many blocks
per setting depends on the tube.} The report claims a feature only at
5$\sigma$ and ranks only settings that reach it. On this instrument
the 90/10 tube gave about 11$\sigma$ per 30\,s block, so two blocks
are plenty; the 1\% \htwo{} Gd-doped tube gave about 2.5$\sigma$ per
block --- two blocks reach 3.5$\sigma$, every setting comes out
``undetermined'', and four reach 5$\sigma$ only just. On a 1\%-class
tube take six blocks per setting (20$k$0 to 20$k$5, about 25\,min for
five settings). Do the normal setting
\emph{within} the ladder, as its own blocks, so the ladder carries
its own anchor even though 2.4 was acquired at the same tuning. For
every setting, \textbf{write down whatever the console shows} for tune
and match --- the readings go into the \texttt{tuning} entries of the
questionnaire (3.2), in whatever units you have. These blocks are
analysed as a ladder of their own and are \textbf{never averaged into
the main noise block} of 2.4. Retune to normal when you are done.

\subsection*{2.7\quad Axis-sign check (${\sim}1$\,min)}

On this console a positive FFT-convention offset points the opposite
way from the physical one: the axis is \emph{mirrored}, the physical
line frequency is carrier $-$ $f_{\mathrm{FFT}}$. That was established
on the partner console two independent ways (vendor checklist item 2):
a $+200$\,Hz \texttt{tof} step moved the apparent line by $+199.6$\,Hz,
and VnmrJ's lock referencing (\texttt{reffrq}, the 0\,ppm frequency)
places water ${\sim}500$\,Hz \emph{below} the carrier where the FFT
convention had it ${\sim}540$\,Hz above. \texttt{tof} itself is
recorded verbatim as \texttt{o1\_hz}, no sign flip. The report still
resolves the sign for every bundle instead of assuming it: from the
one extra acquisition below when it is present, otherwise from the
\file{procpar}'s own referencing (\texttt{reffrq}, \texttt{rfl},
\texttt{rfp} against the water shift), and it flags a disagreement
between the two. The acquisition is the reference of 2.5 once more,
identical in every parameter, with the transmitter offset moved by
exactly $+200$\,Hz. With the probe back at its normal tuning:

\begin{lstlisting}
pw=0.11
tpwr=57
pad=0
at=2
gain=40
tof=tof+200
go
svf('3000_sn_signcal')
tof=tof-200
\end{lstlisting}

The saved \file{procpar} carries the displaced \texttt{tof}; the
converter records it, and the report takes \texttt{tof}(sign check)
$-$ \texttt{tof}(reference) as the displacement. In the physical
convention the water line's apparent offset in the window would
\emph{drop} by exactly that much --- 200\,Hz; on this console it
\emph{rises} by 200\,Hz, which is the mirrored axis, and the report
states which it found. The last line puts \texttt{tof} back so nothing
you run next inherits the shift.

\section{Back at your desk}

\subsection*{3.1\quad Copy the data}

Copy the whole session directory --- every \file{NN\_*.fid} folder ---
onto the laptop, any path. A USB stick is fine; the console never
needs internet access. Copy the console's \file{/vnmr/vnmrrev} (three
lines of text) into the top of that same directory as \file{vnmrrev}.
The packer takes the VnmrJ version from the \texttt{"vnmrj\_version"}
line of \file{answers.json} when that line is present --- the template
below pre-fills SIU's --- and reads this copy only when the line is
left out, so the copy is the fallback that lets you drop the line
rather than retype the version.

\subsection*{3.2\quad Fill in \file{answers.json}}

In the unzipped \file{spin-noise-network} folder, create
\file{answers.json} with the content below --- it is pre-filled for
SIU, so only the \texttt{FILL IN} lines (and \texttt{"locked"},
\texttt{"p90\_us"}, and \texttt{"vt\_setpoint\_k"} if yours differed)
are yours to edit:

\begin{lstlisting}[basicstyle=\ttfamily\footnotesize]
{
  "vendor": "agilent",
  "run_mode": "external-acquisition",
  "facility": {
    "institution": "Southern Illinois University",
    "city": "Carbondale, Illinois",
    "country": "USA",
    "facility_slug": "siu-carbondale",
    "contact_email": "bgoodson@chem.siu.edu",
    "contact_consent": true
  },
  "sample": {
    "description": "FILL IN (e.g. retired lineshape sample, 90% H2O / 10% D2O, Gd-doped)",
    "h2o_fraction_pct": 90,
    "d2o_pct": 10,
    "additives": "FILL IN (dopant, if known)",
    "tube_od_mm": 5,
    "sample_volume_ul": 550,
    "vt_setpoint_k": 298
  },
  "environment": {
    "locked": true,
    "operator_notes": "FILL IN (date; anything unusual during the run)"
  },
  "spectrometer": {
    "probe_type": "RT",
    "console": "Agilent DD2 400",
    "probe_string": "5 mm HCN",
    "coil_temp_k": 296.5,
    "preamp_temp_k": 296.5
  },
  "instrument": {
    "vnmrj_version": "3.2 Revision A",
    "spectrometer_model": "Agilent DD2 400",
    "field_state_notes": "FILL IN (lock on/off during the noise blocks; z0 stable?)"
  },
  "calibration": {
    "p90_us": 10.0,
    "p90_power_db_or_w": "FILL IN (tpwr, dB)",
    "topshim_ok": false
  },
  "experiments": [
    { "expno": 2000, "tuning": { "setting_index": 0, "label": "FILL IN (normal tuning)", "tune_reading": null, "match_reading": null, "units": "FILL IN" } },
    { "expno": 2001, "tuning": { "setting_index": 0, "label": "FILL IN (normal tuning)", "tune_reading": null, "match_reading": null, "units": "FILL IN" } },
    { "expno": 2010, "tuning": { "setting_index": 1, "label": "FILL IN", "tune_reading": null, "match_reading": null, "units": "FILL IN" } },
    { "expno": 2011, "tuning": { "setting_index": 1, "label": "FILL IN", "tune_reading": null, "match_reading": null, "units": "FILL IN" } },
    { "expno": 2020, "tuning": { "setting_index": 2, "label": "FILL IN", "tune_reading": null, "match_reading": null, "units": "FILL IN" } },
    { "expno": 2021, "tuning": { "setting_index": 2, "label": "FILL IN", "tune_reading": null, "match_reading": null, "units": "FILL IN" } },
    { "expno": 2030, "tuning": { "setting_index": 3, "label": "FILL IN", "tune_reading": null, "match_reading": null, "units": "FILL IN" } },
    { "expno": 2031, "tuning": { "setting_index": 3, "label": "FILL IN", "tune_reading": null, "match_reading": null, "units": "FILL IN" } },
    { "expno": 2040, "tuning": { "setting_index": 4, "label": "FILL IN", "tune_reading": null, "match_reading": null, "units": "FILL IN" } },
    { "expno": 2041, "tuning": { "setting_index": 4, "label": "FILL IN", "tune_reading": null, "match_reading": null, "units": "FILL IN" } }
  ]
}
\end{lstlisting}

Three mechanical notes. The two temperatures are the physical lab
ambient you read in 2.1 --- replace 296.5 with your reading, or
\texttt{null} if you have none. Nothing else needs listing: the
converter reads each experiment's role from its save name
(\texttt{\_sn\_ladder}, \texttt{\_sn\_ref\_open}, \texttt{\_sn\_noise},
\texttt{\_sn\_ref\_close}, \texttt{\_sn\_tune}, \texttt{\_sn\_signcal})
and builds the gain ladder from the 27 rungs' stored gains, so the
questionnaire no longer lists noise blocks or ladder rungs. The
\texttt{tuning} entries are the one thing only you know: setting 0 is
your normal tuning, \texttt{tune\_reading} and \texttt{match\_reading}
are the numbers the console showed at that setting (or \texttt{null}),
\texttt{units} names them (``turns'', ``\% reflected'', ``dB''), and
\texttt{label} is free text; add a line per block for settings beyond
4 or for blocks beyond the second (20$k$2 to 20$k$5 on a 1\%-class
tube).

\subsection*{3.3\quad Pack, check, upload}

Three commands, run from inside the unzipped
\file{spin-noise-network} folder:

\begin{lstlisting}
python3 packer/pack_bundle.py --vendor agilent --answers answers.json /path/to/the/data/folder
\end{lstlisting}
Success looks like: the packer validates the session and prints the
path of one \file{spinnoise\_siu-\allowbreak carbondale\_*.zip} ---
and if anything is missing, it lists exactly what, in plain text, so
you fix it and re-run.

\begin{lstlisting}
python3 uploader/upload_bundle.py --selftest spinnoise_*.zip
\end{lstlisting}
Success looks like: the last line reads \texttt{RESULT: PASS} ---
entirely local, no network involved yet.

\begin{lstlisting}
python3 uploader/upload_bundle.py spinnoise_*.zip
\end{lstlisting}
Success looks like: the last line reads \texttt{RECEIPT: <id>} ---
the server has your bundle, checksum-verified on arrival; uploads are
chunked and resumable, so if the connection drops, re-run the same
command and it picks up where it left off.

\section{If anything goes wrong}
\label{sec:wrong}

\begin{tcolorbox}[colback=boxbg, colframe=boxframe,
  title=\textbf{One rule: send it, don't fix it}, fonttitle=\bfseries]
Send the \textbf{exact error text} (copy-paste, not a paraphrase) and
the file it choked on --- the offending \file{NN\_*.fid} folder
zipped, or the bundle zip --- to
\href{mailto:jwbquantum@gmail.com}{jwbquantum@gmail.com}.
The converter has met one real Agilent instrument --- this one, two
sessions --- and every further session is how it becomes trustworthy
beyond that. Nothing you can do here damages anything --- the software
only reads your data, the uploader never deletes your zip, and the
worst case at every step is re-running one command. If the whole chain
fights you, zip the entire session directory and e-mail that (a share
link is fine if it is large); we take it from there.
\end{tcolorbox}

\vspace{0.8em}
\noindent\rule{\textwidth}{0.4pt}\\[2pt]
{\small Contact: John W. Blanchard
(\href{mailto:jwbquantum@gmail.com}{jwbquantum@gmail.com}) ---
CASPEr@night software v0.7.1 --- September 16, 2026. Your raw data remain
yours, and contributing facilities are co-authors on network
publications their data enter.}

\end{document}
